% Compile with XeLaTeX for best Vietnamese support.
% Example: xelatex main.tex

\documentclass[12pt]{article}

\usepackage[a4paper,margin=1in]{geometry}
\usepackage{amsmath, amssymb, amsthm}
\usepackage{hyperref}
\usepackage{xurl}
\usepackage{booktabs}

% XeLaTeX font setup (Windows-friendly)
\usepackage{fontspec}
\setmainfont{Times New Roman}
\setsansfont{Arial}
\setmonofont{Consolas}

\hypersetup{unicode=true, colorlinks=true, linkcolor=blue, urlcolor=blue, citecolor=blue}

% Theorem environments (used by week files)
\theoremstyle{definition}
\newtheorem{definition}{Definition}
\theoremstyle{plain}
\newtheorem{theorem}{Theorem}
\theoremstyle{remark}
\newtheorem*{remark}{Remark}

% -----------------------------
% Report metadata (edit these)
% -----------------------------
\newcommand{\UniversityName}{ĐẠI HỌC QUỐC GIA TP. HỒ CHÍ MINH}
\newcommand{\SchoolName}{TRƯỜNG ĐẠI HỌC KHOA HỌC TỰ NHIÊN}

\newcommand{\CourseName}{Ứng dụng xử lý ảnh số và video số - 23TGMT}
\newcommand{\LecturerName}{Lý Quốc Ngọc - Nguyễn Mạnh Hùng - Phạm Thanh Tùng}

% Put each student on its own line using \\ inside this command.
\newcommand{\StudentList}{%
Lưu Huy Minh Quang - 23127016\\
Trương Quốc Cường - 23127333%
}

\newcommand{\ReportTitle}{BÀI TẬP LỚP NHÀ}
\newcommand{\ReportSubtitle}{Ghi chú / Nội dung học tập}

\begin{document}

\begin{titlepage}
\centering

{\Large \UniversityName\par}
\vspace{0.25cm}
{\Large \SchoolName\par}

\vfill

{\LARGE\bfseries \ReportTitle\par}
\vspace{0.2cm}
{\Large \ReportSubtitle\par}

\vfill

\begin{tabular}{@{}ll@{}}
\textbf{Môn học:} & \CourseName\\
\textbf{Giảng viên:} & \LecturerName\\
\textbf{Sinh viên:} & \begin{tabular}[t]{@{}l@{}}\StudentList\end{tabular}\\
\end{tabular}

\vfill

{\large \today\par}
\end{titlepage}

\tableofcontents
\newpage


\section*{Week 01}

\subsection*{1. Phân biệt các hướng nghiên cứu cơ bản, ứng dụng và phát triển}

Trong \textbf{xử lý ảnh \& video số (Digital Image/Video Processing)}, ba hướng \textbf{cơ bản -- ứng dụng -- phát triển}
có thể hiểu theo lộ trình: \emph{nền tảng lý thuyết} $\rightarrow$ \emph{giải bài toán cụ thể} $\rightarrow$
\emph{triển khai thành hệ thống/sản phẩm chạy ổn định}.

\begin{itemize}
  \item \textbf{(i) Nghiên cứu cơ bản (fundamental/basic).}
  \begin{itemize}
    \item \textbf{Mục tiêu:} xây nền tảng lý thuyết và thuật toán, trả lời ``vì sao'' và ``giới hạn''.
    \item \textbf{Nội dung:} mô hình hoá ảnh/video (nhiễu, mờ, nén, chuyển động), lý thuyết tín hiệu (Fourier/DCT/wavelet),
    tối ưu hoá và bài toán ngược (MAP/ML, regularization), đánh giá chất lượng (PSNR/SSIM/VMAF, tri giác).
    \item \textbf{Đầu ra kỳ vọng:} mô hình/thuật toán mới kèm phân tích (độ đúng, độ ổn định, độ phức tạp), thí nghiệm khoa học
    (ablation, so sánh công bằng, phân tích lỗi).
  \end{itemize}

  \item \textbf{(ii) Nghiên cứu ứng dụng (applied/application-driven).}
  \begin{itemize}
    \item \textbf{Mục tiêu:} giải quyết bài toán thực tế trong một lĩnh vực cụ thể, trả lời ``làm để dùng được''.
    \item \textbf{Nội dung:} phát hiện/nhận dạng/phân đoạn/theo dõi, tăng cường chất lượng ảnh/video, v.v.\ cho các domain như
    y tế (MRI/CT), giám sát an ninh (CCTV), giao thông/xe tự hành, công nghiệp kiểm lỗi, media/phục chế video, viễn thám.
    \item \textbf{Đặc trưng:} dữ liệu thực thường nhiễu và lệch miền (domain shift), thiếu nhãn; KPI theo use-case (mAP/F1, latency,
    robustness); ràng buộc vận hành (thời gian thực, chi phí, quyền riêng tư).
  \end{itemize}

  \item \textbf{(iii) Nghiên cứu phát triển (development/engineering/productization).}
  \begin{itemize}
    \item \textbf{Mục tiêu:} đưa mô hình/thuật toán thành hệ thống triển khai ổn định, tối ưu tốc độ và chi phí.
    \item \textbf{Nội dung:} tối ưu suy luận (quantization, pruning, distillation, TensorRT/ONNX), thiết kế pipeline
    (thu thập $\rightarrow$ gán nhãn $\rightarrow$ huấn luyện $\rightarrow$ kiểm thử $\rightarrow$ triển khai $\rightarrow$ giám sát),
    MLOps/DevOps (versioning, CI/CD, monitoring drift, rollback), tích hợp chuẩn/codec và streaming (H.264/H.265/AV1),
    đảm bảo chất lượng thực tế (robust với ánh sáng/rung/che khuất).
    \item \textbf{Đầu ra kỳ vọng:} prototype/SDK/API hoặc hệ thống realtime với chỉ tiêu rõ (FPS/latency, bộ nhớ, uptime, chi phí).
  \end{itemize}
\end{itemize}


\paragraph{Phân biệt các hướng nghiên cứu}
\begin{itemize}
  \item \textbf{Cơ bản:} tạo \emph{kiến thức/thuật toán nền}.
  \item \textbf{Ứng dụng:} dùng kiến thức đó để \emph{giải bài toán ngành cụ thể}.
  \item \textbf{Phát triển:} biến lời giải thành \emph{hệ thống/sản phẩm chạy ổn, nhanh, rẻ}.
\end{itemize}

\paragraph{So sánh giữa các hướng nghiên cứu}
\begin{table}[h]
\centering
\renewcommand{\arraystretch}{1.25}
\begin{tabular}{p{3.0cm} p{4.2cm} p{4.2cm} p{4.2cm}}
\hline
\textbf{Tiêu chí} & \textbf{Cơ bản} & \textbf{Ứng dụng} & \textbf{Phát triển} \\
\hline
Câu hỏi chính &
``Vì sao / giới hạn / mô hình hoá thế nào?'' &
``Giải bài toán $X$ trong domain $Y$ ra sao?'' &
``Chạy ổn -- nhanh -- rẻ -- dễ vận hành thế nào?'' \\
\hline
Đóng góp kỳ vọng &
Thuật toán/mô hình mới + phân tích (độ đúng, ổn định, độ phức tạp) &
Pipeline + cải thiện KPI cho bài toán thực (mAP, F1, WER, VMAF, \ldots) &
Hệ thống/SDK/production pipeline + SLO (latency, FPS, uptime, cost) \\
\hline
Dữ liệu &
Thường sạch/chuẩn hoá; có thể dùng dữ liệu tổng hợp (synthetic) để kiểm soát biến &
Dữ liệu theo ngành (y tế, giám sát, \ldots), nhiều nhiễu và lệch miền (domain shift) &
Dữ liệu production + logging; theo dõi drift; ràng buộc quyền riêng tư \\
\hline
Đánh giá &
Benchmark + ablation + phân tích lỗi mang tính ``giải thích'' &
KPI theo use-case + độ bền (robustness) + so với baseline thực tế &
Chỉ tiêu hệ thống: FPS/latency, chi phí, ổn định, khả năng mở rộng \\
\hline
Tài nguyên &
Vừa phải $\rightarrow$ nặng (đặc biệt khi huấn luyện mô hình lớn) &
Nặng ở dữ liệu/nhãn và thử nghiệm nhiều điều kiện &
Nặng ở engineering: tối ưu GPU/edge, CI/CD, monitoring \\
\hline
Kỹ năng chủ đạo &
Toán/tín hiệu/tối ưu; thiết kế thí nghiệm chặt &
Hiểu domain + ML/CV + xử lý dữ liệu &
Software/hardware + MLOps + tối ưu hoá và QA \\
\hline
Rủi ro thường gặp &
Ý tưởng hay nhưng khó vượt SOTA/khó thuyết phục &
Chạy tốt trên test nhưng fail ngoài đời (domain shift) &
Mô hình tốt nhưng không chạy kịp/khó vận hành/chi phí cao \\
\hline
\end{tabular}
\end{table}

\paragraph{Yêu cầu của từng hướng nghiên cứu}

\begin{itemize}
  \item \textbf{(A) Cơ bản:}
  \begin{itemize}
    \item Định nghĩa bài toán rõ (noise model, blur model, motion model, \ldots).
    \item Chứng minh/giải thích: vì sao thuật toán hợp lý, trade-off gì, giới hạn gì.
    \item Thí nghiệm ``khoa học'': ablation, sensitivity, phân tích lỗi, so sánh công bằng.
  \end{itemize}

  \item \textbf{(B) Ứng dụng:}
  \begin{itemize}
    \item Use-case và ràng buộc thật (điều kiện ánh sáng, camera, thiếu nhãn, \ldots).
    \item Quy trình dữ liệu: thu thập -- gán nhãn -- chia tập -- chống leakage.
    \item Đánh giá theo KPI đúng nhu cầu (không chỉ PSNR/SSIM; cần tiêu chí nghiệp vụ).
  \end{itemize}

  \item \textbf{(C) Phát triển:}
  \begin{itemize}
    \item Spec hệ thống (FPS/latency, bộ nhớ, throughput, chi phí/1k video, SLA/SLO, \ldots).
    \item Tối ưu triển khai: quantization/pruning, TensorRT/ONNX, batching, streaming pipeline.
    \item Vận hành: versioning dữ liệu/mô hình, giám sát drift, rollback, bảo mật và privacy.
  \end{itemize}
\end{itemize}

\paragraph{Hướng nào đang được đầu tư nhiều nhất hiện nay?}
\begin{itemize}
  \item Nếu xét theo \textbf{dòng tiền của doanh nghiệp/thị trường}, phần \textbf{được rót nhiều nhất} thường là
  \textbf{``Ứng dụng + Phát triển (production hoá)''}---đặc biệt các bài toán \emph{vision AI} gắn với
  sản xuất/giám sát/y tế/giao thông--xe, và triển khai theo mô hình \textbf{edge-to-cloud}.

  \item \textbf{Lý do chính:}
  \begin{itemize}
    \item \textbf{Nhu cầu thị trường tăng mạnh:} nhiều báo cáo dự báo thị trường computer vision tăng nhanh giai đoạn 2025--2030
    (ví dụ các báo cáo thị trường như Grand View Research).
    \item \textbf{Chi tiêu doanh nghiệp dồn vào hạ tầng \& triển khai:} đầu tư cho data center, thiết bị và hạ tầng AI tăng,
    khiến phần ``phát triển/production hoá'' (tối ưu triển khai, vận hành hệ thống) chiếm tỷ trọng ngân sách lớn
    (ví dụ các dự báo chi tiêu CNTT như Gartner).
    \item \textbf{Xu hướng kỹ thuật:} cộng đồng và doanh nghiệp đẩy mạnh hướng \emph{generative + multimodal/foundation models}
    cho thị giác (ảnh/video), kéo theo nhu cầu xây pipeline, công cụ triển khai, tối ưu suy luận (ví dụ các tổng kết nghiên cứu
    tại ICCV 2025 từ Sony AI/Sony).
    \item \textbf{Thực tế triển khai enterprise vision AI:} các báo cáo dựa trên dự án thực tế nhấn mạnh tốc độ đưa mô hình vào
    production, lựa chọn phần cứng, yêu cầu dữ liệu và vận hành (ví dụ Roboflow Vision AI Trends).
  \end{itemize}

  \item \textbf{Tuy vậy,} \textbf{nghiên cứu cơ bản} không ``hết tiền'': các nhóm/lab lớn vẫn đầu tư đáng kể vào
  \textbf{foundation models / generative models} (nền tảng tạo lợi thế dài hạn về thuật toán và dữ liệu), nhưng xét tổng thể,
  ngân sách lớn nhất trong doanh nghiệp thường nằm ở \textbf{hạ tầng + triển khai + sản phẩm} (tức ``phát triển'', rồi đến ``ứng dụng'').
\end{itemize}



\subsection*{2. Cho ví dụ ứng dụng về nghiên cứu cơ bản trong khoa học thị giác}

\subsection*{3. Phân biệt các hướng nghiên cứu cơ bản, ứng dụng và phát triển}

\subsection*{4. Cho ví dụ nhận dạng mẫu sản phẩm và truy vấn mẫu sản phẩm}

\subsection*{5. Nêu ứng dụng cụ thể dùng mô hình học: Supervised Learning, Unsupervised Learning, Reinforcement Learning}

\subsection*{6. Nếu áp dụng meta learning cho Image Classification thì có gì hơn so với Supervised Learning, Objective Detection?}

\end{document}
